\documentclass[11pt]{article}
\usepackage[margin=1in]{geometry}
\usepackage{amsmath,amssymb,amsthm,enumitem,hyperref}
\hypersetup{colorlinks=true,linkcolor=blue,urlcolor=blue,citecolor=blue}
\newtheorem{theorem}{Theorem}
\newcommand{\Res}{\operatorname{Res}}
\title{The asymptotic Genius Conjecture\\\large MathDB 345452}
\author{Principia Math In Progress}
\date{September 18, 2026}
\begin{document}
\maketitle
\begin{abstract}
For every fixed index $i$, we express $F_i-u_i$ as a finite sum of
monomials with rational coefficients in $p$, each of order $O(p^{-1})$.
This proves Federbush's second Genius Conjecture. We also determine the
first correction explicitly.
\end{abstract}

\section*{Task 1. Full solution}
\subsection*{Statement and coefficient formula}
Use the setting of Sections 1--2 of \cite{fed}:
\[
 U(x)=\sum_{k=1}^p u_kx^k,\quad u_1=1,\qquad
 D(x)=\sum_{k=1}^p d_kx^k,\qquad
 A(x)=e^{D(x)}=\sum_{k\ge0}a_kx^k,\quad a_p=0.
\]
The functions $F_1=1,F_2,\ldots,F_p$ are uniquely specified by
\begin{equation}\label{eq:target}
 [x^p]A(x)e^{yU(x)}
       =[t^p]\exp\left(y\sum_{k=1}^pF_kt^k\right).
\end{equation}
Conjecture 2 asks for a finite monomial expansion of $F_i-u_i$ whose
coefficients are rational functions of $p$ tending to zero.
As usual for these indexed rational functions, the limit keeps $i$
fixed and takes integers $p\ge i$ to infinity.

\begin{theorem}\label{thm:main}
For every fixed $i\ge2$ and all integers $p\ge i$,
\begin{align}
 F_i-u_i
 &=\sum_{k=1}^{i-1}\sum_{m=1}^{i-k}
   \sum_{\substack{\ell_1+\cdots+\ell_m=i-k\\\ell_\nu\ge1}}
   \frac{k\prod_{j=1}^{m-1}(i-jp)}
        {m!\prod_{\nu=1}^m(p-\ell_\nu)}
       u_k\prod_{\nu=1}^m a_{\ell_\nu}.\label{eq:finite}
 \intertext{Every rational coefficient in this expression is $O(p^{-1})$.
After expanding the $a_\ell$ into monomials in the $d_k$, this proves
Conjecture 2. More precisely, coefficientwise in the variables,}
 F_i&=u_i+\frac1p\sum_{k=1}^{i-1}k u_kd_{i-k}
                   +O(p^{-2}).\label{eq:first}
\end{align}
\end{theorem}

\subsection*{A formal coordinate formula}
For a self-contained proof, we recall the coordinate construction
used for the companion linearity result \cite{companion}. Put
\[
 W_p(x)=\sum_{k=0}^{p-1}\frac{p}{p-k}a_kx^k,\qquad
 \phi(x)=xW_p(x)^{-1/p},\qquad \psi=\phi^{-1}.
\]
Fractional powers have constant term $1$, and $\phi=x+O(x^2)$
has a unique formal compositional inverse. Differentiation gives
\[
 \frac{\phi'}{\phi^{p+1}}
   =x^{-p-1}\left(W_p-\frac{x}{p}W_p'\right)
   =x^{-p-1}\sum_{k=0}^{p-1}a_kx^k.
\]
Since $a_p=0$, formal residue substitution yields, for every
power series $H$,
\[
 [x^p]A(x)H(x)
   =\Res_x\frac{H(x)\phi'(x)}{\phi(x)^{p+1}}
   =[t^p]H(\psi(t)).
\]
Apply this to $H=e^{yU}$. It follows that $F_i=[t^i]U(\psi(t))$.
Uniqueness can also be seen directly: the coefficient of
$y^{p-i+1}$ in \eqref{eq:target} determines $F_i$ with multiplier
$1/(p-i)!$ once the lower $F_k$ have been determined.
Formal integration by parts now gives
\begin{equation}\label{eq:lagrange}
 F_i=\Res_x\frac{U(x)\phi'(x)}{\phi(x)^{i+1}}
   =\frac1i[x^{i-1}]U'(x)W_p(x)^{i/p}
   =\sum_{k=1}^i\frac{k}{i}u_k[x^{i-k}]W_p(x)^{i/p}.
\end{equation}
These are standard formal residue and Lagrange--B\"urmann operations;
no analytic limit or convergence assumption is involved.

\subsection*{Finite rational coefficients and their limits}
In \eqref{eq:lagrange}, the term $k=i$ is exactly $u_i$.
For $k<i$, expand the binomial series
\[
 W_p^{i/p}=\sum_{m\ge0}\binom{i/p}{m}
       \left(\sum_{\ell=1}^{p-1}\frac{p}{p-\ell}a_\ell x^\ell\right)^m.
\]
Only $1\le m\le i-k$ and positive compositions
$\ell_1+\cdots+\ell_m=i-k$ contribute. Their scalar coefficient is
\[
 \frac{k}{i}\binom{i/p}{m}
        \prod_{\nu=1}^m\frac{p}{p-\ell_\nu}
 =\frac{k\prod_{j=1}^{m-1}(i-jp)}
        {m!\prod_{\nu=1}^m(p-\ell_\nu)},
\]
proving \eqref{eq:finite}. Its numerator has degree at most $m-1$
in $p$ and its denominator has degree $m$. Also
$\ell_\nu\le i-1<p$, so no denominator vanishes in the allowed range.
Thus every displayed coefficient is a rational function that is
$O(p^{-1})$.

For each $\ell\le i-1$, the coefficient $a_\ell$ is the finite
polynomial
\begin{equation}\label{eq:monomials}
 a_\ell=\sum_{\substack{b_1,\ldots,b_\ell\ge0\\
                         b_1+2b_2+\cdots+\ell b_\ell=\ell}}
           \prod_{s=1}^\ell\frac{d_s^{b_s}}{b_s!}.
\end{equation}
Substituting these polynomials into \eqref{eq:finite} gives a finite
list of monomials $u_k\prod d_s^{b_s}$. The list depends on $i$,
but not on $p\ge i$. Each coefficient is a rational constant times
one of the rational functions above, or a finite sum of them, so
each tends to zero. This is precisely the asserted monomial
decomposition. Notice that $d_p$ does not occur: the constraint
$a_p=0$ therefore imposes no extra restriction on the lower variables
used in this formula.

\subsection*{The first correction}
For a fixed composition in \eqref{eq:finite}, its scalar coefficient
has expansion
\[
 \frac{k\prod_{j=1}^{m-1}(i-jp)}
        {m!\prod_{\nu=1}^m(p-\ell_\nu)}
       =\frac{(-1)^{m-1}k}{mp}+O(p^{-2}).
\]
For $s=i-k\ge1$, the finite sum of the leading terms in the
$a_\ell$ is
\[
 \sum_{m=1}^s\frac{(-1)^{m-1}}m
  \sum_{\substack{\ell_1+\cdots+\ell_m=s\\\ell_\nu\ge1}}
     \prod_{\nu=1}^m a_{\ell_\nu}
   =[x^s]\log A(x)=d_s.
\]
This proves \eqref{eq:first} and completes the theorem.\hfill$\square$

\subsection*{Verification and scope}
Four exact verification groups check the original identity, finite
rational formula, formal expansion through second order, and explicit
$F_2,F_3$ formulas. Their 255 cases use exact fractions and supplement
the proof for arbitrary fixed $i$.

All six pages of the primary were read. Conjecture 2 is distinct from
the linearity statement in Conjecture 1, treated in \cite{companion}.
The setting is Federbush's; the Bell formulation is credited there to
David Williams. Bounded checks found no journal final or prior external
full target claim as of September 18, 2026. Review is by the authoring
agent only, without independent peer review or proof-assistant
certification. The limit is coefficientwise at fixed $i$.

\section*{Tasks 2--5. Proposed follow-up directions}
\begin{enumerate}[label=\textbf{Task \arabic*.},start=2,leftmargin=*,itemsep=2pt,topsep=4pt]
\item Derive a general closed recurrence for all coefficients of the
fixed-index expansion of $F_i$ in inverse powers of $p$.
\item Establish useful estimates when $i$ also grows with $p$, including
the regime $i/p\to\alpha\in(0,1)$.
\item Classify cancellations among the rational coefficients in
\eqref{eq:finite}, and determine minimal common denominators.
\item Apply the exact transformation to the auxiliary coefficient
corrections in the graph-positivity setting, with explicit dependence
on the degree and perturbation parameters.
\end{enumerate}
\begin{thebibliography}{9}
\bibitem{fed} P. Federbush, \emph{The Genius Conjectures (via Bell
Polynomials)}, \href{https://arxiv.org/abs/2002.03814v1}{arXiv:2002.03814v1},
2020. Conjecture 2, page 2.
\bibitem{companion} Principia Math In Progress,
\emph{Linearity in Federbush's Genius Conjecture},
\href{https://github.com/antoshashakov/Principia-Math-In-Progress/tree/mathdb-results/mathdb-open-problems/problems/345451}{MathDB 345451 solution package},
September 18, 2026. The coordinate argument is reproduced above.
\end{thebibliography}
\end{document}
